% ================================================================
%  第 3 章 一封信 (后半: 原书 p.92-108, 边码 L31-L56)
%  源文: _work/src1/task-ch03b.md
% ================================================================
%
% ---------------- p.92 ----------------
\origpage{L31}现成的观念;关于我自己,关于一个我从前与之认同的圈子,关于一些我爱过、也许如今依然爱着的人——与这些观念“分离”,或者更确切地说,让它们从我身上自行脱落。这是一件工作,确确实实是工作,但从来不是一场搏斗——这工作带给我的,除了许多别的可贵之物,还有时而出现的悲伤时刻,却从来没有一刻悔恨,也没有一刻苦涩。苦涩是回避一种认识的手段之一,是回避一段经历所传消息的手段之一;是以关于世界、关于他人的另一种“幻觉”(不妨说是负片式的)为代价,把自己维持在对自身的一种顽固幻觉之中。

\origpage{L32}我是没有苦涩、也没有悔恨地看着这些现成的观念一件一件从我身上脱落,像卸下一件件累赘、甚至压人的重负——它们曾为我所“珍爱”,出于积习,也因为它们“从来”就在那里。它们确已变成了某种第二本性。但这“第二本性”并不是“我”。把它们一块一块地剥离,并不是撕裂,甚至也不是那种被剥夺了珍视之物的人的失落。“剥离”一词,在我这里是作为工作的报偿与果实而来的。它的标志,是一种当下即至的、有益身心的轻松,一种受欢迎的解放。

\section{(10) 运动中的四道浪}

果不其然,这封信与我动笔时所设想的全然不同。我本来主要是想就“葬礼”做一份小小的“概述”:大致情形如此如此,你信也罢不信也罢(我自己当初也几乎难以相信……),可事情就是这样,无可置疑,不管你乐意不乐意,某年某月某日某页,白纸黑字印在某期刊某本书上,去看就是了——再说,《收获与播种》里一切都已逐条拧开来细讲;见“四种运算”诸注——爱信不信!你要是宁可不去读我,自有别人替你去读……

结果这些全都没有。可这封信已经写到三十页的坎上了,而我原计划统共五六页。也不是我存心如此:一页一页写下来,我被引着对你讲的,竟是那些最本质的东西;而我当初那么急不可耐要倒空的“口袋”(这一回它可真是明明白白摆在头几页上!),到现在还没解开呢!它甚至不再挠我的手指了,那点欲望在路上就消散了。我明白了:这里不是地方……

说实话,《收获与播种》的第四部分(也是最长的一部分),名为“葬礼(3)”或“四种运算”,原本出自一则“按语”,当初正是打算做成一份小小的“概述”,好把去年那场突袭式(而且风风火火)的调查所揭示的东西大致勾勒一下——那场调查续在第二部分里(“葬礼(1)”,或“中国皇帝的袍子”)。我以为那不过是一则五页或十页的“按语”,不会更多。结果一环扣一环,调查又被它带着重新跑了起来,写成了将近四百页——几乎是我本该为之做摘要、做总结的那部分的两倍!于是那份小小的概述至今仍然缺着,而《收获与播种》里已有六百页是专门调查“葬礼”的。这确实有点蠢。不过,在把稿子交给印刷厂之前,总还来得及把它补进《导论》的第三部分(那部分早就不是在乎多十页二十页的时候了)。

《收获与播种》的五个部分(最后一部分尚未写完,恐怕几个月之内也写不完),是(三道)“冥想”之浪与(两道)“调查”之浪的交替。这里仿佛以一种缩影的方式,映照出我最近九年的生命——那九年同样也是“浪”的交替,那些浪涌自如今支配我生命的两种激情:冥想的激情与数学的激情。而说实话,\origpage{L33}%
% ---------------- p.93 ----------------
我方才用“调查”这个信手拈来的名字所称的那两个部分(或“两道浪”),恰恰是直接从我作为数学家的过往的扎根处涌出来的,被我对数学的激情、以及扎根于这激情之中的种种自我执念所驱动。

第一道浪,“愚妄与更新”,是与我的数学家过往的第一次相遇,由此引出对我当下的一次冥想——我刚刚发现,这当下正扎根于那段过往之中。诚然,这一点绝无半点预谋;可这一部分为《收获与播种》此后全部篇幅定下了“基调”,它像是一种内在的、天赐而不可或缺的准备,好让我承受紧随其后的“光华四射的葬礼”之发现——那是在第二道浪里,“葬礼(1)——或中国皇帝的袍子”。

其实,与其说这是一场“调查”,不如说这就是那场发现逐日展开的故事:它对我是怎样的冲击,我如何竭力面对这样不打招呼就劈头砸下来的东西,如何设法用我的亲身经历去安放那难以置信之物,直到它终于变得为我所熟悉,好歹让它变得可以理解。这一运动引出一个最初的、暂定的终点,在按语《掘墓人——或整个教团》(第97号)里:这是第一次尝试,要在一个多年来、如今更比任何时候都更尖锐地显出“对常理的可怕挑战”之面貌的东西里,辨认出一种解释、一种意义!

这同一场第二运动,也引出一次“病中插曲”\fnote{36}{}{这次插曲是两则按语的主题:《事件——或身体与精神》和《陷阱——或轻易与耗竭》(第98、99号),它们开启了名为《亡者(始终没有故去)》的“第十一支队列”。},逼我绝对静养,三个多月里中止了一切智力活动。那正是我自以为又一次即将把《收获与播种》收尾的时候(只差最后几件“后勤”活儿……)。去年九月底前后恢复常态活动,我正准备终于给那些还陷在困境里的按语做最后的收尾,仍以为只消再添两三则收尾的按语,其中一则讲我刚经历的那次“健康事件”。结果,一周又一周、一月又一月,又来了上千页——比已经写好的还多出一倍——而这一回,显然我还是没写完\fnote{37}{}{“始终没有写完”——哪怕只因为后面还应有一道第五部分,而在写这几行的时候,它还没有完成。}!

事实上,在这段漫长的中断期间,我几乎与一种(在搁下它时)再热不过(甚至滚烫!)的质料失去了接触;这几乎强迫我以新的眼睛回到那质料上去——如果我不想愚蠢地“了结”一个“程序”的最后收尾,而那个程序我早已与它断了活的联系。

第三道浪就是这样诞生的,诞生在《收获与播种》这一片浩瀚的运动之中——那是一道漫长的“冥想之浪”,主题是阴与阳,是事物的动态与人存在中“暗”与“明”的两面。这道冥想出自一种愿望:更深入地理解在“葬礼”中起作用的那些深层力量;然而它从一开始就获得了自身的自治与统一,并且一上来就朝向最普遍的东西,同时也朝向最私密个人的东西。正是在这道冥想之中,我发现了这件事(说实话它显而易见,只要人肯把问题提出来):在我自发地发现事物的进程里,无论在数学还是在别处,“基调”都是“阴”的、“女性”的;而且尤其重要的是,与最常见的情形相反,我始终忠实于我身上这种原初本性\fnote{38}{}{不过,这种“对我原初本性的忠实”绝非彻底的。长久以来,它仅限于我的数学工作;而在别的一切地方,尤其在我与他人的关系里,我随大流,抬高并优先看待我身上那些被感受为“男性”的特征,压抑那些“女性”的特征。关于这一点,在《一个生命的故事:三乐章的循环》(第107—110号)那一组按语里有相当详细的讨论,那组按语差不多是《阴阳之钥》的开篇。},从未让它弯折或修正,以
%
% ---------------- p.94 ----------------
\origpage{L34}迁就周遭圈子里被尊奉的主导价值。这一发现起初在我眼里只是一种有趣的好奇。只是渐渐地,它才显出是理解“葬礼”的一把关键钥匙。而且——这一点在我看来意义还要大——如今我极其清楚、不留丝毫疑问地看到:如果说,以绝非出众的智力天赋,我竟然能始终在数学工作中使出自己的全部本事,创造出一部作品,孕育出一种广大、有力而丰饶的视野,那么这完全要归功于这种忠实,归功于这种全不理会去迁就规范的心态——靠了它,我才怀着完全的信任委身于原初的求知冲动,不去削它、砍它,不减去任何造就其力量、其精细、其不可分割之本性的东西。

\origpage{L35}然而,在这道名为“葬礼(2)——或阴阳之钥”的冥想里,处在注意力中心的并不是创造力及其源头,而毋宁是“冲突”,是创造力的阻塞状态,是创造的能量因心灵中敌对力量(多半是隐秘的力量)的对抗而涣散。暴力的那些面貌——那种(表面上)“无端的”、“为了好玩”的暴力——曾不止一次让我在“葬礼”中不知所措,并让一大堆相似的亲身处境重新浮现。对这种暴力的经验,在我一生中就像是“冲突经验那坚硬的、不可化约的核”。我此前从未直面过这样一个可怕的谜:这种暴力在一般人类存在中、尤其在我自己的存在中,其存在本身与普遍性究竟是怎么回事。整个下半部(即“阴”的一面,或“衰”的一面),注意力都集中在这个谜上。正是在冥想的这一部分里,对“葬礼”之意义、对其中表达出的那些力量,一种更深的看法逐渐浮现出来。这大概也是《收获与播种》在自我认识层面上最丰饶的一部分:它让我接触到一些要害的问题与处境,并且恰恰让我感觉到那种“要害”的性质——而直到去年,这一点还一直被回避着。

在这段关于阴阳的没完没了的“离题”走完之后,我发觉自己大体上仍旧剩着那“两三则按语”要写(顶多加一两则,其中一则早已有了现成的名字:“四种运算”……),写完了就与《收获与播种》两清。

后来发生的事大家都知道了:这“最后几则按语”最终成了《收获与播种》最长的一部分,将近五百页。这就是这场运动的“第四道浪”。它也是“葬礼”的第三即最后一部分,我给它取名“四种运算”——这也是那一组按语(“对一具尸体做的四种运算”)的名字,那组按语构成了这第四次思考吐息的核心。在《收获与播种》里,这一部分才是严格意义上的“调查”部分——不过还带点调味:这场调查并不局限于纯粹的“技术”方面、即“侦探”方面;这里的思考,如同《收获与播种》别处一样,首先是由求知与求理解的欲望驱动的。这里的调子当然比“葬礼”第一部分要“带劲”得多——那会儿我还有点揉着眼睛、自问是不是在做梦呢!话虽如此,一页页披露出来的事实却常常来得正是时候,让许多此前在这里那里只是一带而过的东西当场得到说明,落实到具体而触目的例子上。也正是在这一部分里,数学上的离题占了重要位置,那是被一次重新接上的接触激发起来的(出于调查的需要)——与一种我十五年未见的质料重新接触。在光谱的另一端,还有对吾友左格曼·梅布胡特(Zoghman Mebkhout)遭遇的现场叙述(这一部分即献给他):他落到一伙手段高强、毫无顾忌的“黑手党”手里,而当初他投身于那个(固然引人入胜、表面看来又无伤大雅的)课题——各类簇的上同调——时,根本没有梦到过这些。为了给
%
% ---------------- p.95 ----------------
\origpage{L36}这整个“调查”部分那纠缠错综的按语、副按语、副副按语……的迷宫提供一条简明的线索,我请你去看目录(第167$'$号至1767号),以及这一叠里的第一则:《侦探——或玫瑰色的人生》(第167$'$号)。不过我要提一句:这则按语注明4月22日,后来有点被“事态甩在了后面”——因为一波三折,这场我当时以为(基本上)已经查到底的调查,又拖拖拉拉地继续了两个月。

\origpage{L37}这第四次吐息一口气延续了四个多月,从二月中旬直到六月底。主要是在思考的这一部分里,靠一种细致而顽强的“档案工作”,日子一天天过去、页面一页页累积,“葬礼”的现实才逐渐建立起一种具体的、可触的接触;总之,我才好歹与它“熟识”起来——尽管它在我身上引起的(并且继续引起的)那种发自脏腑的拒绝反应,一直阻碍着真正的认识。这段漫长的思考,以对德利涅来访的回顾开篇(本信中已经提到过),而以关于我与塞尔的关系、关于塞尔在“葬礼”中的角色的一番“最后一刻”的思考收尾\fnote{39}{}{见按语《家庭相册》(第173号)的 c、d、e 三部分,其中最后一部分注明的日期是6月18日(正好是十天前)。只有一则按语或按语片段日期更晚(即《一场屠杀中的五篇论文——或孝道》,第1767号,注明次日,6月19日)。你会注意到,在《收获与播种》这第四部分、即“调查部分”里,与别处相反,按语往往按逻辑顺序而非时间顺序相随。例如“葬礼”最后两则按语(构成结尾的《自深处》)注明4月7日,比我方才引的那则按语早两个半月。我还要指出:除了《葬礼(3)》真正的“调查”部分(第167$'$—1767号按语)——那部分构成殡葬大典的“第五个时辰”(《阴阳之钥》是第二个时辰)——按语都按写作顺序相随,极少例外。}。此前那个空缺,就在于我暗中把塞尔“排除在外”,借着我已说过的那个“禁忌”——如今我觉得,这或许正是直到上个月为止我对“葬礼”的理解中留下的最严重的空缺;于是这番“最后一刻”的思考,反倒成了《收获与播种》这“第四次吐息”带给我的最重要的东西:让我对“葬礼”及其间表达出的力量,有一种不那么单薄、更加充实的领会。

\section{(11) 运动与结构}

我想,关于《收获与播种》我有什么最要紧的话想对你说,都已兜了一圈说完了;为了让你先知道“这是怎么回事”,我说的肯定已经足够,足够让你判断:接下去那一千多页的信,是否“与你有关”——因而,你是否还要继续读下去。若答案是“是”,那么我觉得,再就《收获与播种》的形式补几句说明(尤其是实用方面的)是有用的。

这形式是某种精神的映照与表达,我竭力想把那种精神在前面的篇幅里“传递”出来。与我过去的出版物相比,如果说《收获与播种》里(以及它所出自的《追寻栈》里)出现了一种新的品质,那大概就是自发性。当然,这里有贯穿的线索,有大的追问,它们赋予整个思考以连贯与统一。然而思考是逐日进行的,没有预先立定的“纲领”或“计划”,从来没有什么“应该证明的东西”事先定给我。我的意旨不是证明,而是发现,是往一种未知的质料里更深地走进去,是让那些还只是被预感、被猜疑、被瞥见的东西凝结起来。我可以毫不夸张地说,在这件工作里,没有一天、也没有一夜的思考是在“预料之中”的范围内展开的——无论就观念、意象、联想而言,那些东西都是我坐到
%
% ---------------- p.96 ----------------
\origpage{L38}白纸面前、要在上面顽强地追一条线索、或者接上另一条刚冒出来的线索时,已经在那里的。每一次,思考中出现的东西都与我事先能够预言的不同——假使我竟敢勉力去预先描述我以为在眼前看见的东西的话。最常见的是,思考走上了开头全然未料的道路,通向同样未料的新风景。然而即便它沿着一条大致预定的路线走,几小时里旅程向我显示的东西,也与我上路时对它抱有的形象相去之远,正如一处真实的风景——带着它清凉的阴影与温暖的光的嬉戏,随着散步者脚步而微妙变换的透视,还有被一阵吹得草舞、吹得芦苇歌唱的微风送来的无数声响与无名芬芳……——这样一处活的、不可捉摸的风景,与一张明信片相去之远,无论那明信片多么美、多么成功、多么“准确”。

是那种一口气追下去的思考——一天之内、或一夜之内——构成不可分割的单位,可以说是活的、个体的细胞;在整片思考(《收获与播种》在此即是)之中。整场反思之于每一个这样的单位(或那些共同构成旋律的“按语”\fnote{40}{}{最初写《愚妄与更新》时,“按语”这个名字对我就是“注解”的同义词,起的是脚注的作用。出于排版上的方便,我宁愿把这些注解放到正文末尾(第1至44号按语,第141页与第171页)。这样做的原因之一是:这些“按语”或“注解”中有一些要占上一页乃至好几页,甚至比它们本该评论的正文还要长。至于思考“初稿”那些不可分割的“单位”,我当时因为找不到更好的名字,就把它们叫作“节”(比“段落”不那么令人生畏!)。\par 这一情形和文本的结构,随着下一部分而改变;那一部分最初叫“葬礼”,后来成了“葬礼(1)”(或“中国皇帝的袍子”)。这段思考接在双重按语《我的孤儿们》与《拒受一份遗产——或一个矛盾的代价》(第46、47号,第177、192页)之后,那两则按语是给《收获与播种》最后那一“节”(更确切地说,是给它本来要有的第一部分、即《愚妄与更新》)《一段过去的重量》(第50号,第131页)作注解的。后来,又给同一节添了别的注解(第44$'$、50号按语),还有别的按语来给《我的孤儿们》作注解,这些又依次生出新的作注解的按语;这一回还不算上真正的脚注——当预定的注解(写成白纸黑字之后仍然是)篇幅不大的时候。于是,从理论上说,《收获与播种》那一部分(当时它本该构成第二即最后一部分)就表现为一组给《一段过去的重量》那一“节”作的“按语”。由于惯性,这种“按语”(而非“节”)的细分法在随后三个部分里仍旧保留着;在那里,我给思考“初稿”作注解的手段兼用两种:脚注(当篇幅允许时),以及正文中引到的、写在后面的按语。\par 从排版上说,“按语”与“节”(在 RS I 里用作思考“初稿”的基本单位)的区别在于一个记号,如(1)、(2)等(即把按语编号放在括号里、排成上标,这是引注的通行做法),它或者放在该按语的开头,或者作为一种引注放在正文中提及它的适当位置。而“节”则用1至50的阿拉伯数字标出(不用那些令人生畏的下标与上标——出于实际需要,我在按语里不得不使用它们)。话虽如此,可以说《收获与播种》第一部分里“节”的功能与后面各部分里“按语”的功能并无本质区别。我在本信这一部分(《自发与结构》)中关于这一功能所作的评论,同样适用于 RS I 的“节”,尽管我一律用“按语”这个通名。\par 至于别的细则与约定,尤其是关于如何读《葬礼(1)》的目录,我请参见《导论》第7节(《葬礼的编排》),特别是第 xiv–xv 页。}……),正如一个活机体的身体之于它每一个个体细胞——那些细胞千差万别,每一个都占着一个只属于它的位置、担着一个只属于它的功能。不过有时候,在同一段一口气追下来的思考里,事后会觉察到一些重要的断口,把它分成好几个这样的单位或消息,于是每一个都得到自己的名字,由此获得自己的身份与自治。另一些时候则相反,一段由于某种原因(多半是偶然原因)被打断的思考,第二天或第三天又自发地接了下去;或者一段横跨两三天写成的思考,事后看来却像是从头到尾一口气写成的;仿佛只有睡眠的需要,才迫使我们违心地在那里插入某种断口(可以说是“生理性的”),它只由一行干巴巴的日期
%
% ---------------- p.97 ----------------
\origpage{L39}(有时还由好几行)标出来,标在所谈那则“按语”的某些相邻段落之间——这时那则按语便以一个独一无二的名字把自己区分出来。

\origpage{L40}就这样,《收获与播种》的每一则按语都有它自己的个体性,有使它区别于任何别的按语的面孔与功能。对每一则,我都设法用一个名字表达它独有的特性;那名字应当还原或唤起它“要说”的东西之本质,至少是某种本质。每一则,我真正认它,首先就是凭它的名字;此后每当需要它来帮忙,我也是凭这个名字唤它。

名字往往在我还没想到它之前就自己来了。正是它不期而至的出现,向我示意:我还在写的这则按语已近完成——它已经把要说的话说完,只等我写完眼下这一段……同样常见的是,在继续思考之前重读昨天或前天的按语时,名字也照样自发地出现。有时,在刚来的那则按语出现之后的几天或几周里,名字会有所变动,或者添上第二个我起初没想到的名字。许多按语有两个名字,表达对它所传消息的两种不同的照明,有时是互为补充的照明。这些双名中第一个向我呈现的,是在“愚妄与更新”开头不久,即《与克劳德·谢瓦莱相遇——或自由与善意》(第11号)。

只有两次,我在开始写一则按语之前就先有了名字——而且两次都后来被事态的发展搅乱了!

只有拉开距离——几周、甚至几月的距离——之后,在逐日相继的按语整体里,才会显出某种整体的运动与结构。我设法靠各种分组与再分组来把握这两者,每一组都有它自己的名字,那名字赋予它独有的存在、它的功能或消息;有点像方才那个比喻里同一个身体的器官与四肢,以及四肢的某些部位。于是,在《收获与播种》这“整体”里,有我已说过的五个“部分”,每一部分各有自己的结构:“愚妄与更新”分组成八“章”,I 至 VIII\fnote{41}{}{在《愚妄与更新》里,我有时把这些章称作《收获与播种》的“部分”;当然不要把它们与已说过的那五个部分相混,那五个部分是后来才出现的。};而构成“葬礼”的那三个部分合起来(它们也是几个月里逐渐显出来的……),则由一支漫长而庄严的、十二支“队列”I 至 XII 的行列组成。其中最后一支,更确切地说是那场“殡葬大典”(这就是它的名字)——前面十一支“队列”都朝它走去(肯定没怎么料到什么……)——其规模真是巨大,与它所隆重安葬的那部作品的规模相称:它囊括了 RS III(《葬礼(2)》)的几乎全部和 RS IV(《葬礼(3)》)的全部,连同它近八百页、一百五十来则按语(而当初,这场大名鼎鼎的典礼只预定包含两则!)。这场典礼由大主祭亲自(以他众所周知的谦逊……)手法娴熟地主持,分为九个彼此分开的“时辰”或礼仪环节,以《悼词》开场(这还用说),并以《自深处》收尾(理当如此)。这些“时辰”之中另有两个,一名《阴阳之钥》,一名《四种运算》,每一个都(远远)构成它所嵌入的《收获与播种》那部分(III 或 IV)的最大份额,并且顺带把自己的名字给了那一部分。

在《收获与播种》全书中,我像护着眼珠一样照料目录,不断改动它,以顾及源源不断涌来的未预料的按语\fnote{42}{}{这些未预料的按语里,尤其有那些“出自一则尺寸变得无法接受的脚注”的。最常见的是,我把它紧放在它所隶属的那则按语之后,给它同一个编号,必要时加上一个 $'$ 或 $''$ 号,甚至加上 $'''$——这就免去了把此后所有已写好的按语一并重新编号这桩干不下来的活儿!这些由一则脚注生出另一则的按语,在目录里前面带一个 $!$ 号(至少在《葬礼(1)》里如此)。},并
%
% ---------------- p.98 ----------------
\origpage{L41}且尽我所能让它精细地映照出思考的整体运动与其中显现的微妙结构。正是在第三、尤其是第四部分(刚才说到的“钥匙”与“四种运算”)里,这结构最为复杂、最为交错。

\origpage{L42}为了给文本保住自发性的品格、保住思考如其实际进行与实际被经历时那种未预料的面貌,我不愿意把名字放在每则按语前面——何况名字每次都是事后才出现的。所以我建议你:每读完一则按语,就翻到目录去看看这则按语叫什么;顺便也可以一眼看出它如何嵌入已经进行的思考,甚至嵌入尚未来到的思考。否则,你就有望无可救药地迷失在一堆看来难以消化、五花八门的按语里,其编号有时古怪,甚至可说令人生畏\fnote{43}{}{至于这些有时看来也许古怪的编号何以必要,我请你参见这封滔滔不绝的信里前一条脚注。};就像一个旅人在一座陌生的城市里迷了路(那城市是按世代与世纪的任性古怪地长出来的……),既没有向导,连一张地图也没有,帮他辨认方向。在准备付印的稿子里,我打算随文加入“章”以及别的按语组与节组的名字,只把按语(或节)本身除外。但即便那样,偶尔查阅目录在我看来仍不可少,免得迷失在几百则按语的一团乱麻里——它们一串接一串,绵延一千多页……

\section{(12) 自发与严格}

自发与严格,是同一品质那“暗”与“明”的两面。唯有二者的婚配,才生出文本或存在者的那种特殊品质——人们也许可以用“真之品质”这样的说法去唤它。如果说在我过去的出版物里,自发性(即便不是缺席,也是)被压缩到极小的份额,我并不认为,随着它在晚年于我身上绽放,严格性因此就变小了。毋宁说,它那位阴性的伴侣的完全在场,给了严格性一种新的维度、新的丰饶。

这种严格也针对它自己施行:它留意使自己必须做的那桩微妙的“筛选”——在意识领域里流过的大量东西中,不断地把有意义的、本质的东西从偶然的、次要的东西里澄清出来——不至于变厚、凝固成审查与自满的自动机制。唯有我们心中求知的渴望,才能唤醒并激发这样一种不笨重的警觉、这样一种灵敏,去对抗那巨大而无处不在的惰性,即那些由现成观念凿成的“(所谓)天然的坡道”——它们是我们的恐惧与我们所受的塑造的表达。

而这同一种严格、同一种警觉的注意,也朝向自发性,也朝向那些带自发性面貌的东西,为的是在那里同样分辨出成分:分辨出那些固然再自然不过的“坡道”,把它们与真正从存在深处、从原初的求知与行动冲动中喷涌出来的东西区别开来——那冲动把我们带向与世界的相遇。

在写作层面上,严格表现为一种持久的用心:用语言尽可能精细、尽可能忠实地圈定要表达的思想、情感、知觉、意象、直觉……;在所要表达之物轮廓分明之处,不满足于一个含糊或近似的词;而在所要表达之物仍被“还只是被预感”的迷雾围绕之处,也不满足于一个虚假精确(因而同样走形)的词。当我们试图就地捕捉它如其所是的样子时——也只有在那时——未知之物才向我们显露其真正的本性,也许甚至直到大白天光之中,如果它生来属于白昼、而我们的欲望又促它卸下阴影与雾的帷幕。我们的角色不是妄称描述并固定我们所不知、从我们手中逃逸的东西,而是谦卑地、炽烈地去认识那从四面八方围绕着我们的未知与神秘。

% ---------------- p.99 ----------------
\origpage{L43}这就是说,写作的角色不是记录一项研究的结果,而正是记录研究的过程本身——记录爱的工作,记录我们与我们的母亲、那“世界”、那“未知者”的爱的作品;她不知疲倦地召唤我们进入她之中,去认识她那取之不尽的躯体里的她,去到她之中那神秘的欲望之路把我们带去的任何地方。

为了传达这一过程,那些回头的修正——它们给写作的“初稿”、甚至第二稿第三稿添上细微差别、加以精确化、深化,有时加以改正——本身就是发现进程的一部分。它们构成文本本质的一部分,赋予文本全部意义。所以,放在《愚妄与更新》末尾的那些“按语”(或“注解”)——构成文本“初稿”的那五十“节”里处处引到它们——是文本不可分割的、本质的一部分。我竭力劝你边读边翻过去看,至少每读完一节的末尾(那节里出现了一处或几处对这类“按语”的指引)就去看。对《收获与播种》其余部分里的脚注也是如此;或者,当某则“按语”(在此充当“正文”)里引到后面某一则按语,那后一则便起“回看”它或注解它的作用。这一点,连同我请你不要在阅读中途离开目录的劝告,就是我要给你的阅读建议中最主要的一条。

最后一个问题,是实用性的,它将(有点平淡地)结束这封该收尾的信。有些时候有点“慌乱”,为了把《收获与播种》的各个分册准备好,交给系里的复印部去印,好让印刷(如果可能)赶在大假之前。匆忙之中,有一整页最后一刻的脚注要补进第2分册(《葬礼(1)——或中国皇帝的袍子》),却“跳”掉了。那些主要是改正某些直到最近、在写《四种运算》的过程中才出现的材料性错误。这些脚注里有一条比其他几条更重要,我想在此指出。那是对按语《受害者——或两种沉默》(第78$'$号,第304页)的一则注解。在那则按语里,我竭力——除了别的——去圈定我的一些印象(当然全是主观的):关于吾友左格曼·梅布胡特当时如何“内化”他所遭受的那种令人不安的掠夺;而这则按语被他觉得对他不公,因为我似乎几乎把他与那些掠夺者“装进了同一个袋子”。可以肯定的是,在那则按语里——它并不自称给出别的东西,只给出与某个特定“时刻”相关的印象——我只敲响了一口钟,而把另一些同样真实(也许还不那么可争议)的钟声留在未言之中(而且大概还当作不言自明的事)。不管怎样,对这一个微妙题目的思考,在一年之后,在按语《根与孤独》(第171号)里大大深化了。这一则没有引起左格曼的保留。关于同一题目的另一些思考材料,也见于《三块界石——或无辜》与《死页》(第171(x)与(xii)号)这两则按语。这三则按语都属于《封神》,那是《四种运算》里专讲对左格曼·梅布胡特作品的据为己有与窃取那场运算的部分。

我只剩下祝你阅读愉快了——也盼着轮到我读你的信!

亚历山大·格罗滕迪克

\origpage{L44}
\subsection*{后记:或一场辩论的语境与前情}
\tocsubsub{后记:或一场辩论的语境与前情}

% ---------------- p.100 ----------------
\origpage{L45}
\section{(13) 瓶装光谱仪}

这封信写完之后,整整七个月过去了;寄出去也已近四个月,连同随信而去的那个“大砖头”。而且每一本上都有我手写的题赠\fnote{44}{}{只有少数例外,主要是那些我本人不认识的同行,他们只收到试印本的第0和第4分册,作为他们积极参与我的“葬礼”的奖赏。}。像一只“漂流瓶”,更确切地说,像一整群这样的漂流瓶子,我的消息一路漂去、流转,直到这个我曾熟悉的数学微观宇宙最偏远的角落。随着日子、星期、月份一天天过去,直接与间接的回声传回来,于我不期然地,面前仿佛出现了一张数学界的巨大射线照片——那是由一架触手四伸的“瓶装光谱仪”拍下的,而我那些天真的“瓶子”就是它一根根游走的天线。于是(责任随身份而来!),我这个本来并不缺事可做的人,如今面前又多了一件新任务:破译这张射线照片,并且尽我所能,把我从中读出的东西报告出来。这将作为《收获与播种》的第六部分(也是最后一部分,我保证!)。倘若上帝假我以岁月,这一部分便将为我“晚年那部社会学巨著”加冕。眼下,先说几句初步的评论。

用来迎接我这支可怜的手工小舰队的,占压倒优势的,是一种甜腻又带三分凶气的调子,唱的是“瞧,格罗滕迪克老了,犯起偏执狂来啦”,或者“这位可真把自己当回事”——如此这般,戏就演完了!然而这种调子的信我只收到一封\fnote{45}{}{这封信出自一位曾是我的学生的人,而且他还是我的一位“同葬者”。},另有两封属于那种含蓄的、自得其乐的嘲弄\fnote{46}{}{出自两位当年在布尔巴基共事的同行,其中一位还是我初出道时善意接待过我的年长者之一。}。我的数学家收信人,包括那些曾是我的学生的人,大多数以沉默作答\fnote{47}{}{寄给数学家的共一百三十一份,至今有五十三个收信人有了动静,哪怕只是告知收到。这其中有六个是我过去的学生——另外八个人里,我一个也没得到音信。}——这沉默对我说了很多。

不过我已经收到了大量通信。绝大多数信是客气而尴尬的调子,那尴尬常常还想显得友好,仿佛是出于礼数。有两三次,我在这尴尬背后、仿佛被它滤过之后,感到一种依然活着的情感的温度。最常见的情形是:尴尬若不是以对好情感的郑重声明来表达(为他自己,或为别人),就是以恭维来表达——我这一辈子也没收到过这么多恭维!唱的是“伟大的数学家”、“了不起的篇章”(关于创造力“以及诸如此类”……)、“无可争议的作家”,等等等等。为了凑足分量,我甚至还有幸得到一句情真意切(绝非反讽)的恭维,夸我内心生活之丰富。不用说,在所有这些信里,我的通信人都小心翼翼地不进入任何问题的实处,更不会亲自卷入;调子倒更像是那种“被请来发表意见”的人(借用其中一封信的说法)面对一桩有点不干不净的事——何况还是假设的、甚至想象的——总之尤其要紧的是:一桩与他本人无关的事。可是当他装出要碰一碰其中一个问题的样子时,那是指尖轻触,并且尽可能把它推得离自己远远的——或者借给我一番好劝告之机,或者用谨慎的条件式,或者在不知\origpage{L46}%
% ---------------- p.101 ----------------
该说什么时用那套现成的套话,或者用别的什么法子。\origpage{L47}不过还是有些人透露说,也许发生过一些不太正常的事——同时小心地把是什么事、是什么人留在最大的含糊之中……

我也收到过坦率热烈的回声,来自十五六位新老朋友。有些人表达了感动,并不想掩饰它、也不想让它噤声。这些回声,以及另一些来自数学界之外、同样热烈的回声,将是我对一项漫长而孤独的工作的报偿——那工作不仅为我自己,也是为所有人做的。

在收到我这封信的一百三十来位同行中,有三个人真正在完整意义上作了回应:他们把自己放了进去,而不是限于对时下事件作一番远距离的评论。我还从一位非数学界的通信者那里收到一个这样的回声。这些才是对我的消息的真正答复。这也正是我最好的报偿。

\section{(14) 三只脚踩一个盘子}

我的好几位同行和数学家朋友曾表示希望,《收获与播种》能在数学界引发一场广泛的辩论:关于这个圈子里的风气状况,关于数学家的伦理,关于他工作的意义与目的。目前,至少可以说,事情并没有往那个方向走。从现在起(说一句应景的双关语),关于一场“葬礼”的辩论,看样子已经径自被一场辩论的“葬礼”取代了!

不管人们愿不愿意,尽管大多数人的沉默与冷漠,一场辩论却确确实实已经开了场。它恐怕永远不会具有真正公开辩论那样的规模,更不会(上帝保佑别这样!)具有“官方”辩论那种排场与僵硬。无论如何,已经有很多人抢先一步、三下两下就把它了结了——在还没弄清它是什么之前,就在自己心里把它关掉了;他们仗着那条万古不易、不可动摇的共识:“在这最好的世界里,一切都再好不过”(在此即数学的世界)。不过,也许终究会有一次质疑从外面来,渐渐地从一些“证人”那里来——他们不属于同一个圈子,因而不是这个圈子的群体共识的囚徒,因而(即便在自己心里)也不觉得自己被牵连在内。

在收到的几乎全部回声里,我都看到同一种混淆,涉及两个前提性的问题:《收获与播种》所提出(至少是默契地提出)的“辩论”究竟针对什么;以及谁有资格了解它并就它表态,或者说:谁有资格在充分知情的基础上形成意见。关于这一点,我想在此标明三个“坐标点”。当然,这并不能阻止那些执意要混淆的人继续待在混淆里。至少,对于那些想知道究竟是怎么回事的人,这也许能帮他们不被四面八方来的响动(甚至包括那些最善意的响动……)分了神。

a)一些真诚的朋友向我保证“一切最终都会摆平”(而“一切”,我猜,指的是那些不幸被弄坏了的“东西”……);我只要重返舞台,“用新工作让自己站住脚”,做做演讲等等——其余的别人会做。人们会宽宏地说“对这个该死的格罗滕迪克,咱们也的确有点不公”,然后不声不响地、带着或多或少的诚意去拨正航向\fnote{48}{(*)}{我已经不止一次注意到这类微妙的迹象,表明人们已经好好记下:狮子醒了……};甚至慈父般地拍拍他的肩膀,说他是“伟大的数学家”,好让这么个毕竟还算可敬的家伙平静下来——可惜他偏偏作势要发火,要掀起不受欢迎的风浪。

正如这些朋友所暗示的,这里根本不是要“松松绳子”、也不是要让人松松绳子。就我而言,我全然不需要恭维,不需要真诚的仰慕者,也不需要“盟友”——为“我的”事业,或为任何别的事业。问题不在我身上——我过得好得很;也不在我的作品——它自己会说话,哪怕是对聋子说。如果说这场辩论也(在别的之外)涉及我这个人、我的作品,那只是因为它们是对别的东西的揭示——通过一场葬礼的现实(的确是最具揭示性的现实之一)。

如果真有“某个人”在我看来该唤起警觉、不安与紧迫之感,那绝不是我这个人物,甚至也不是我的任何一位“同葬者”。那是一个集体的存在者,既不可捉摸又十分具体,人们常常谈起它,却从来小心翼翼地不去审视它——它的名字叫“数学共同体”。

最近这几周,我终于把它看成一个有血有肉的人,而它的身体正遭受着深处的坏疽。最好的食物、最精美的菜肴,在它体内都变成了毒,使病害传播得更广、扎得更深。然而它又有一种不可遏制的贪食,要永远吃得更多,仿佛这肯定是它自欺的一种方式——关于一种它无论如何也不愿去认识的病。对它说什么都是白费——最最简单的词也失去了意义。它们不再是消息的载体,只剩下引发恐惧与拒绝的喀嗒声……

b)我的大多数同行,甚至那些心肠好的老朋友,在贸然发表意见时,都要围上一圈谨慎的条件式,诸如“如果事情真是……那确实是不能容许的”——好让他们心满意足地再躺回去睡大觉。可我原以为我说得够清楚了……

\origpage{L48}隔着七个月的距离,如今我可以说明:对于《收获与播种》里报告并评论的几乎全部事实,它们的真实性不成为任何争论的对象。我会在后面回到那少数几个例外(它们也会各自在相应的地方被标明是这样的);至于其余全部事实,在《收获与播种》初版写成之后,与几位主要当事人的仔细对证(即皮埃尔·德利涅、让–皮埃尔·塞尔与吕克·吕西耶)消除了细节上的错误,并就材料性事实本身达成了毫不含糊的一致\fnote{49}{(*)}{我很高兴向这三位表达我的感谢,感谢他们在此事上表现出的善意;并确认他们在一切涉及材料性事实的问题上完全诚实。}。

所以,辩论根本不涉及事实的真实性——那不成问题——而涉及这样一个问题:这些事实所描述的作派与态度,是否应当被视为被认可的、“正常的”,还是不是。

这里说的是这样一些作派:在我的见证里,我(也许错了……)称之为丑闻式的;如背信弃义或滥用权力,如明目张胆的不诚实,不止一次达到了不义与无耻的程度。在了解了这些事实(十五年前还不可想象的事实)之后,我还得学的一件事——相当难以想象的一件事——就是:我的同行中的大多数,甚至包括那些曾是我的学生或朋友的人,今天竟把这些作派视为正常,视为完全体面。

c)我的许多同行与老朋友维持混淆还有第二种办法。调子是这样的:“抱歉,可我们不是这一行的专家——别要求我们去了解那些(天佑如此……)从我们头顶上飞过去的事实……”

我则断言:要了解那些主要事实,根本不必是“专家”(这回轮到我抱歉了!),甚至不必会背乘法口诀表或知道毕达哥拉斯定理。

% ---------------- p.103 ----------------
\origpage{L49}连《熙德》或拉封丹《寓言》都没读过也不要紧。一个发育正常的十岁孩子,在这件事上跟最负盛名的专家一样有能力(甚至比他还强……)\fnote{50}{(**)}{当然,我写《收获与播种》并不是给十岁孩子看的;若要对他说话,我会选一种他熟悉的语言。}。

\origpage{L50}让我举一个例子来说明这一点,就从“葬礼”里随手拈来一个\fnote{51}{(***)}{这是我在1984年4月19日发现的第一场“葬礼大手术”,也正是那一天,“葬礼”这个名字强加给了我。关于此事,见同一天写的那两则按语:《一个梦的回忆——或母题的诞生》与《葬礼——或新父亲》(RS III,第51、52号)。那里也有即将谈到的这本书的完整出处。}。要了解下面这几件事并就它们作出判断,根本不必知道“Motif”这个多面而极微妙的数学概念的来龙去脉,也不必小学毕业。

1°) 1963年至1969年间,我引入了“Motif”这一概念;并围绕这一概念发展出一套“哲学”与一套“理论”,它们部分地仍停留在猜想阶段。无论对错(此处无关紧要),我把 Motif 理论看作我为我那个时代的数学带来的最深的东西。再说,如今“Motif 的瑜伽”的重要性与深度已无人质疑(在它被近乎完全沉默了十年之后——那沉默从我离开数学舞台之后随即开始)。

2°)在唯一一本(1981年出版)、基本上以 Motif 理论为专题的书里(这个由我引入的名字就出现在书名中),唯一一处能让读者疑心我这点微末之人与书中长篇大论发展的那种理论有远近关系的段落,在第261页。那段话(两行半)是在向读者解释:那里发展的理论,与一个名叫格罗滕迪克的人的理论毫无关系(该理论在那里被提到第一次也是最后一次,没有任何别的出处或说明)。

3°)有一个著名的猜想,人称“霍奇猜想”(它究竟讲什么都无所谓),它的成立将意味着那本光彩夺目的书里所谓“另一种”Motif 理论,与我将近二十年前在众目睽睽之下发展出的那种理论相同(是它的一个极特殊的特例)。

我还可以加上 4°):那本书四位署名者中最负盛名的一位曾是我的学生,他多年来那些精彩的想法,除了从我这里,还能从哪里学来?——而他在书里把它们呈现得像是刚刚才想出来似的\fnote{52}{(*)}{我并不是说这本书里没有一些想法,甚至一些漂亮的想法,出自这位作者或别的合著者。但整本书的问题框架、赋予它意义的概念语境,乃至那套微妙的 X-范畴理论(被误称为“Tannaka 式”)——那套技术构成了全书的核心——这些都是我的成果。};5°)这两件事在消息灵通的人中间是众所周知的,可是在文献里白费力气也找不到一处书面痕迹,能证明那位出色的作者也许从我口中听到过什么\fnote{53}{(*)}{不过有一个例外:1977年塞尔写的一份报告里有一行,将在相应之处谈到。};以及 6°)那个微妙的算术问题——据主要作者本人向我解释,它构成全书的核心问题(而我的名字始终未被提起)——是我在六十年代沿着“Motif 的瑜伽”的路子提出来并澄清的,而作者是通过我才知道它的;我还可以再堆上 7°、8° 等等(在相应之处我当然不会不这么做)。

前面这些,对我的用意来说已经够了,我的用意就是:要了解这类事实并就它们作出判断,根本不需要什么专门的“能力”——“事情不在那个层面上”。这里起作用的能力,除了健全的理性(原则上人人都有)之外,就是我愿意称之为体面感的东西。

这里说的那本书,如今已是数学文献中被引用最多的著作之一,而它的“主要作者”也是当代最负盛名的数学家之一。话虽如此,而且看得清楚,在这桩事情里,如今在我看来最最值得注意的是:在这本书数不清的读者当中,包括那些第一手知道究竟是怎么回事的人,包括那些曾是我的学生、我的朋友的人——竟没有一个人看出其中有什么不正常。无论如何,直到我写这几行的今天,还没有一个人来向我表示,对这本光彩夺目的书有丝毫保留\fnote{54}{(**)}{总共只有两位同行(包括左格曼·梅布胡特)向我表示过这样的“保留”。这两位都不能算是这本书的“读者”。他们是出于好奇翻了翻,想弄弄清楚……}。

% ---------------- p.104 ----------------
\origpage{L51}至于我的同行和老朋友中那些从未把这书拿在手里、却以此为由申辩自己无能的人,我要对他们说:要当“专家”,根本不必——只要到随便哪一家数学图书馆去把那卷书借来,翻一翻,亲眼看看那无人争议的东西就行了……

\section{(15) 坏疽——或时代精神(1)}

\origpage{L52}这场“Motif 运算”只是同一水准的四场“大运算”之一,此外还有一大批规模较小、精神相同的。它远不是用来充实我那份“一个时代的风俗画”的集体蒙骗中“最大”的一桩,更不是最独特的一桩。它所做的,不过是趁富人不在(或已故……)之机,只掠夺了他的羊群;而不是(在普遍的漠然中)为了取乐,当着他的面去掐死穷人的母羊。甚至直到数学语言里如今已进入日常用法的层面,那些看来无害的书名、概念名、命题名,时时被人引用,它们本身就已经是一种蒙骗或欺诈\fnote{55}{}{我这里想到的尤其是那个古怪的缩写“SGA 4 ½”(分数还是有用处的!),它一个就抵得上双重欺诈(而且是当代数学文献中被引用最多的缩写之一);还有“维尔迪耶对偶”或“维尔迪耶对偶对象”、“德利涅–格罗滕迪克猜想”之类的名称,最后还有“Tannaka 式范畴”(这一回 Tannaka 本人并无过错,因为从来没人问过他……)。这些将在相应之处更详细地谈到。},并以各自的方式见证一个时代的不体面。

如果我自认为曾为“数学共同体”做过什么有用的事,那就是把若干不光彩的事实带到了大白天光之下——那些事实原本在暗处发酵。这肯定是人人都天天、或差不多天天、远远近近接触到的那类事实。他们之中有多少人肯费点闲工夫停下来——哪怕只是一瞬——去闻一闻空气,看一看?

那个自己曾撞上一些人的傲慢、另一些人(或同一批人)的不诚实的人,也许还自以为那是一种专门分配给他的特别霉运。把自己的经验与我的见证相对照,也许他会感到,这“霉运”也是他给一种时代精神起的名字,那种时代精神压着他,一如压着所有人。而且(谁知道呢!)也许这会促使他投入一场辩论——那场辩论与他有关,正如与我有关。

可是,如果我“拿到公共广场上摊开”的这堆“脏衣服”,除了引来一些人没有笑意的冷笑、另一些人客气的尴尬,在众人的漠然中别无反响,那么一个本来浑浊的局面就会变得非常清楚。(至少对那个还愿意用自己眼睛的人是这样。)真诚与体面的传统共识\fnote{56}{}{我说这些“出于真诚与体面的共识”,并不是说它们从未被违反过。但即便在被违反的时候,那也确实是“违反”,而共识本身仍然是被接受的。},在数学家之间的关系里、在数学家与其技艺的关系里,从此都将成为过去之物,“过时”之物。无需有哪个国际数学家组织再隆重宣告,这从此也将是既成之事、几乎是官方之事:如今,一切手段都允许,不再有任何保留与限制——对于数学界里掌权的那伙“靠内部推举结成的兄弟会”而言。一切观念上的捣鬼
%
% ---------------- p.105 ----------------
\origpage{L53}把戏,只为把那个只等着相信的漠然读者牵着鼻子走;一切父权名分的黑市交易;同伙之间的假引注,以及对注定沉默者的沉默;各种拉帮结派与形形色色的伪造,乃至在众目睽睽之下最粗劣的剽窃——是的,对这一切都说“阿门”,并且带着数学公共广场上所有“大名头”、所有大大小小老板的祝福,无论是用言语还是用沉默(若不用积极而殷勤的参与的话)。是的,对那里风靡一时的“新风格”说“阿门”!那曾经是一门艺术的东西,如今经(几乎)一致的赞同,变成了一个搅浑水、抢地盘的集市,在老板们慈父般的目光下。

曾经有一个时候,在数学家的世界里,权力的行使受到一致的、不可触犯的共识的限制,那是集体体面感的表达。这些共识与这种感情,如今都成了过时之物、被超越之物,肯定配不上计算机、太空舱与中子弹的光荣时代。

从此这将是既成而封定之事:权力,对于那伙掌握它的人而言,是一种任意专断的权力。

\section{(16) 公开认错——或时代精神(2)}

在《信》里,我想我已充分说明了我写《收获与播种》时所怀的精神,足以让人清楚:我绝不声称在那里做史家的工作。那是一份真诚的见证,关乎第一手的亲身经历,以及对此经历的思考。见证与思考都供所有人取用,包括史家——他可以把它当作众多材料之一来用。那时,把这份材料交付批判性分析、依其技艺的严格准则来做,便是他的事。

当然,应当区分狭义上的事实(“生事实”或“材料性事实”)与对这些事实的“评估”或“解释”——后者赋予它们意义,而这意义对一位观察者(或共同行动者)和对另一个人并不相同。大体上可以说,《收获与播种》的“见证”方面关涉事实,而它的“思考”方面关涉对事实的解释,也就是我为赋予它们意义所做的工作。在构成见证的“事实”里,我也归入“心理事实”,尤其是我的见证所映照的各种情感、联想与意象,无论它们发生在或远或近的过去,还是就在写作的当下。

对于我在《收获与播种》里描述或提及的事实,我区分三类来源。有些事实是记忆还给我的,其清晰或模糊程度因场合而异,有时还被变形。关于这些,我可以为我在写作之时的真诚立意担保,却绝不为全无错误担保。相反,我曾有机会指出其中若干细节错误,并在相应之处用后来补写的脚注标明。另一方面,有书面文献,尤其是书信,特别是正式发表的科学出版物,我有时以所期望的全部精确性引到它们。最后,还有第三方的证词。有时它补充我自己的记忆,让我重新唤起、加以精确化,有时加以改正。在某些少见的场合(我马上会回过来谈),这证词给我带来了与我已知信息全然不同的新信息。当我转述这样一份证词时,这并不意味着我有能力自始至终核查它的准确与可靠,而只是说:它相当可信地嵌入了我第一手所知的那张丰富的事实之网,足以让我确信(无论对错……)这份证词在本质上确实合乎真相。

% ---------------- p.106 ----------------
\origpage{L54}对于一个细心的读者,我想在任何时候都不会有困难去“分清”事实的记述与对事实的解释,并且(在前一种情况下)辨明我方才描述的三类来源中,是哪一类在起作用。

\origpage{L55}\asterisms

方才我提到一份第三方证词,我转述了它,却未能“自始至终核查其可靠性”——那就是左格曼·梅布胡特的证词,关于围绕他著作的那场大规模掉包手术。在《收获与播种》里提到的“材料性事实”中,如今尚可争议、或按我现在的判断需要改正的,只有某些仅凭梅布胡特一人证词而立的事实。为结束这篇后记,我要在此对试印本《收获与播种》中呈现的“梅布胡特事件版本”提出一些批评性的评论。更详细的评论与更正将各自在相应之处收入印刷版(即《收获与播种》的定本)。

我曾想充当其转述者的那个“梅布胡特版本”,在我看来本质上由以下两个论点构成:

1. 1972年至1979年间,梅布胡特是孤身一人\fnote{57}{}{1975年柏原的可构造性定理是例外,它在理论中的重要性毫无争议。但按梅布胡特的说法,那该是柏原对这一正在诞生的理论的唯一贡献。这一(不准确的)说法,曾由柏原其他论文的缺席所佐证——在那些论文里,他本该至少暗示某些主要想法。},在普遍的漠然中,从我的著作汲取灵感,发展出“$D$-模哲学”,把它作为我意义上的“上同调系数”的新理论。

2. 无论在法国还是在国际上,都有过一种一致的共识,要把他的名字与角色从这一新理论中抹掉——一旦它的意义开始被承认。

这个版本有很强的文献支持:一方面是梅布胡特那些完全令人信服的出版物;另一方面是其他作者许多出版物的署名(尤其是1981年6月吕米尼讨论会论文集),其中掉包的蓄意用心不容任何怀疑。最后,梅布胡特另外提供给我的那些更详细的细节(我在“葬礼(3)——或四种运算”那部分里转述了它们),虽然无法直接核实,却与某种整体的氛围完全吻合,而那氛围的真实性在我看来已不容任何怀疑。

我刚得知若干新的事实\fnote{58}{}{我感谢皮埃尔·沙皮拉与克里斯蒂安·乌泽尔,他们愿意提醒我注意这些事实,以及我对梅布胡特–柏原之争的陈述中那种偏颇之处。},它们表明对上面第 1° 点须作大大的修正。梅布胡特所处的孤立\fnote{59}{}{这种孤立首先来自我过去的学生对梅布胡特的想法与工作的漠然——他固执地装作在追随一位被一致共识判给遗忘的“先祖”……}确实是真的,但那是一种相对的孤立。在法国,有 J.-P. 拉米在同一个题目上的工作(梅布胡特对我只字未提),而尤其是:看来梅布胡特所发展并完成到底、且被他归为自己名下的某些重要想法,可能应归于柏原\fnote{60}{}{这些想法中最重要的是所谓“黎曼–希尔伯特”对应(用新风格的术语来说)在 $D$-模上的版本。相关猜想已由梅布胡特证明,并且(据沙皮拉向我肯定)也由柏原证明(而梅布胡特向我保证他的证明是唯一发表的)。关于证明优先权的问题,对我仍是一片迷雾,我放弃把余生用来把它弄清楚……\par 至于以 $\nabla^{\infty}$-模表述的姊妹命题,看来毫无疑问:想法的归属与证明都属于梅布胡特。}。这一来,就使我(过于)忠实地转述的那个梅布胡特版本里关于柏原–梅布胡特之争的某些情节,显得不可信或可疑了。

% ---------------- p.107 ----------------
\origpage{L56}毫无疑问,无论就“档案工作”的层面而言,还是就他设法完成到底的某些想法的构想而言,梅布胡特都是 $D$-模新理论的主要拓荒者之一,也许甚至是首要的拓荒者;无论如何,他是唯一一个全身心投入这项任务的人,而这项任务的真正意义当时还没向他显现,正如它也没向所有人显现。也确实,围绕这部著作发生的那场掉包手术——以吕米尼讨论会为顶点——在我看来仍是数学界这个世纪的大不体面之一。但若声称(像我出于真诚所做的那样)梅布胡特是独自担此任务的人,那就错了。相反,他是唯一一个有诚实与勇气明确指出我的想法与著作在他的工作中、在新理论的萌发中所占分量的人。

这篇后记不是细谈此事的地方——我会在相应之处细谈,包括一些有助于说明“梅布胡特版本”心理语境的评论。如果说“梅布胡特–柏原的纠纷”对我有某种意义,那仅仅在于它照亮了一个时代的整体氛围。而对我来说,甚至它那些变形、甚至通过那些促其浮现而起作用的力量,“梅布胡特版本”也显得是我为“一个时代的卷宗”提供的众多不那么可争议的材料之中一份有力的“时代征兆”。

我还须为我的轻率作公开认错:我呈现梅布胡特–柏原之争时,只用了一幅仅仅依据梅布胡特所提供的证词与文献的图画,而且仿佛那个版本不容任何怀疑。那个版本把一位第三方置于可笑乃至可憎的光线下,这就更有理由谨慎行事。为我的轻率、为这种缺乏健全的谨慎,我很愿意在此向柏原先生表示我最诚挚的歉意。
